\begin{abstract}
As generative artificial intelligence has become increasingly sophisticated, the ability to effortlessly produce hyper-realistic manipulated media---often referred to as deepfakes---has introduced profound challenges to verifying digital authenticity. For a long time, traditional deepfake detection methods have heavily relied on training massive neural networks to identify spatial and semantic inconsistencies directly within raw RGB images. While effective against early generation deepfakes, these semantic cues are becoming vanishingly rare as generative models improve their visual blending capabilities. Furthermore, these massive networks carry extreme computational overhead. 

In this work, we pursued an alternative, fundamentally different approach to digital forensics: rather than looking at what the image portrays, we examined the digital fingerprint of how the image was saved. We employed Error Level Analysis (ELA) to expose the underlying compression artifacts. Because manipulated or spliced regions typically exhibit a different JPEG compression history compared to the original background, ELA highlights these discrepancies as stark, high-frequency visual noise. We processed a 10,000-frame subset of the standard FaceForensics++ dataset, converting natural images into ELA heatmaps. Recognizing the domain shift between natural imagery and ELA noise, we fine-tuned an EfficientNet-B0 architecture using a strategy of differential learning rates. By completely unfreezing the pre-trained base and allowing the lower-level convolutional filters to adapt to compression grids, our optimized model achieved a robust validation accuracy of 86.40\% and an F1-score of 86.11\%. Ultimately, we demonstrate that compression artifact patterns serve as highly reliable and lightweight discriminatory features for deepfake detection.
\end{abstract}
